% 03_trace_subspace_formulation.tex
%
% Purpose:
%   Present the main mathematical trace-subspace formulation, including lead
%   Bloch modes, incoming/outgoing Cauchy subspaces, the DtN interpretation, and
%   the center-cell port matching matrix.
%
% Main algorithm:
%   Build lead-cell scattering matrices, solve the Bloch generalized
%   eigenproblem, select admissible Cauchy trace subspaces, and couple them to
%   the center Muller boundary integral row.
%
% Based on:
%   sections/bloch_leads.tex, sections/dtn_relation.tex, and
%   sections/bie_center.tex from the first report draft.
%
% Main changes:
%   Combines the conceptual core of the method into one paper-style top-level
%   section with several subsections.
%
% Numerical goal:
%   Record the block matrix and dimension count used by the scattering and
%   eigenvalue workflows.

\section{Trace-Subspace Formulation}

\subsection{Lead-cell scattering and Bloch multipliers}

For a single periodic lead cell with left wall \(x=X_L\) and right wall
\(x=X_R\), write local Rayleigh amplitudes as \(a^L,b^L,a^R,b^R\), where
\(a\) denotes a local right-going Rayleigh amplitude and \(b\) denotes a local
left-going Rayleigh amplitude.  The finite-dimensional cell scattering relation
is
\[
  \begin{bmatrix}
  b^L\\
  a^R
  \end{bmatrix}
  =
  S_{\mathrm{cell}}
  \begin{bmatrix}
  a^L\\
  b^R
  \end{bmatrix}
  =
  \begin{bmatrix}
  R_L & T_{R\to L}\\
  T_{L\to R} & R_R
  \end{bmatrix}
  \begin{bmatrix}
  a^L\\
  b^R
  \end{bmatrix}.
\]
The Bloch condition for a multiplier \(\lambda\) is
\[
  a^R=\lambda a^L,
  \qquad
  b^R=\lambda b^L .
\]
With \(z=[a;b]\), this gives the generalized eigenproblem
\[
  \Asc z=\lambda\Bsc z,
\]
where
\[
  \Asc=
  \begin{bmatrix}
  -R_L&I\\
  T_{L\to R}&0
  \end{bmatrix},
  \qquad
  \Bsc=
  \begin{bmatrix}
  0&T_{R\to L}\\
  I&-R_R
  \end{bmatrix}.
\]

\subsection{Incoming and outgoing trace subspaces}

The Bloch eigenvectors are converted to Cauchy traces on lead-cell walls.  On a
cell left wall,
\[
  D_L(:,j)=a_j^L+b_j^L,
  \qquad
  N_L(:,j)=\ii\GammaR(a_j^L-b_j^L),
\]
where
\[
  \GammaR=\diag(\gamma_m).
\]
The right wall traces satisfy
\[
  D_R(:,j)=\lambda_jD_L(:,j),
  \qquad
  N_R(:,j)=\lambda_jN_L(:,j).
\]
Here \(N_L,N_R\) denote \(\partial_x u\) trace coefficients on a lead-cell wall.
They are converted to center-cell outward-normal derivatives only when attached
to \(\Gamma_-\) or \(\Gamma_+\).

For the positive half-lead, outgoing decay as \(x\to+\infty\) corresponds, away
from unit-circle complications, to
\[
  |\lambda_j|<1 .
\]
Thus
\[
  D_{\mathrm{out}}^+=D_L(:,\mathcal I^+),
  \qquad
  N_{\mathrm{out}}^+=N_L(:,\mathcal I^+),
  \qquad
  \mathcal I^+=\{j:|\lambda_j|<1\}.
\]
For the negative half-lead, outgoing decay as \(x\to-\infty\) corresponds to
\[
  |\lambda_j|>1 .
\]
Since the outward normal at the center-cell negative port is \(-e_x\),
\[
  D_{\mathrm{out}}^-=D_R(:,\mathcal I^-),
  \qquad
  N_{\mathrm{out}}^-=-N_R(:,\mathcal I^-),
  \qquad
  \mathcal I^-=\{j:|\lambda_j|>1\}.
\]

The incoming and outgoing trace subspaces are
\[
  E_{\mathrm{in/out}}^\pm
  =
  \operatorname{range}
  \begin{bmatrix}
  D_{\mathrm{in/out}}^\pm\\
  N_{\mathrm{in/out}}^\pm
  \end{bmatrix}
  \subset \CC^K\times\CC^K .
\]
Rayleigh coordinates are not Bloch modes.  They provide a wall coordinate basis,
while the Bloch computation identifies the physical admissible subspaces.  Any
well-conditioned basis of the same selected subspace may be used numerically,
for example after QR or SVD.  An arbitrary subspace of \(\CC^{2K}\), however,
would impose an artificial boundary condition.

\subsection{Relation to DtN maps}

At the continuous level, a positive half-lead outgoing Dirichlet-to-Neumann map,
when it is well defined, has the form
\[
  \Lambda_{\mathrm{out}}^+:
  H^{1/2}(\Gamma_+)\to H^{-1/2}(\Gamma_+),
  \qquad
  D^+\mapsto N^+ .
\]
An analogous map may be written for the negative half-lead.  The exterior PDE
alone does not define this map; one must impose a radiation, admissibility, or
limiting absorption condition to select the exterior branch.

The implementation keeps the Cauchy trace condition in subspace form:
\[
  \begin{bmatrix}
  D^\pm\\
  N^\pm
  \end{bmatrix}
  \in
  E_{\mathrm{out}}^\pm .
\]
If the selected subspace is a graph over its Dirichlet projection, then
\[
  N^\pm=\Lambda_{\mathrm{out}}^\pm D^\pm .
\]
Near resonances, thresholds, or points where this projection loses rank, forming
an explicit \(\Lambda_{\mathrm{out}}^\pm\) can be ill-conditioned.  The
Cauchy-subspace form avoids this unnecessary graph requirement and matches the
data naturally produced by the Bloch eigenproblem.

For forced problems, uniqueness is expected only away from resonances and
eigenvalues.  For homogeneous eigenvalue problems, a nontrivial admissible
nullspace is the target.  Limiting absorption, for example replacing \(k\) by
\(k+\ii\varepsilon\) with \(\varepsilon>0\), can regularize mode selection and
half-guide solution branches, but it should not be confused with changing the
final lossless eigenvalue problem.

\subsection{Center-cell trace matching formulation}

In the center background region, the homogeneous Rayleigh part is written
\[
  u_h(x,y)
  =
  \sum_{m=-M}^{M}
  a_m\ee^{\ii\gamma_m(x-X_-)}\psi_m(y)
  +
  \sum_{m=-M}^{M}
  b_m\ee^{-\ii\gamma_m(x-X_-)}\psi_m(y).
\]
The vectors \(a,b\in\CC^K\) are local center-cell right-going and left-going
Rayleigh coefficients.  They are not positive and negative lead indices.

Let
\[
  E=\diag(\ee^{\ii\gamma_mL_0}),
  \qquad
  L_0=X_+-X_-,
  \qquad
  \GammaR=\diag(\gamma_m).
\]
The center-cell outward-normal traces of \(u_h\) are
\[
  D_h^- = a+b,
  \qquad
  N_h^- = -\ii\GammaR(a-b),
\]
and
\[
  D_h^+ = Ea+E^{-1}b,
  \qquad
  N_h^+ = \ii\GammaR(Ea-E^{-1}b).
\]

The exterior center-cell field and interior obstacle field are represented as
\[
  u_e=u_h+u_s^e,
  \qquad
  u_i=u_s^i .
\]
With the density convention
\[
  \eta=
  \begin{bmatrix}
  \tau\\
  -\sigma
  \end{bmatrix},
\]
the center obstacle transmission condition gives the forced Muller row
\[
  \Aqp\eta+H_\Sigma^a a+H_\Sigma^b b=0 .
\]
The matrices \(H_\Sigma^a,H_\Sigma^b\) are the Cauchy data on
\(\Sigma=\partial\Omega_0\) of the center homogeneous Rayleigh basis.

The center density produces outgoing Rayleigh coefficients at the ports:
\[
  s^- = F_-\eta,
  \qquad
  s^+ = F_+\eta .
\]
Thus
\[
  D_s^- = F_-\eta,
  \qquad
  N_s^- = \ii\GammaR F_-\eta,
\]
and
\[
  D_s^+ = F_+\eta,
  \qquad
  N_s^+ = \ii\GammaR F_+\eta .
\]
The same \(+\ii\GammaR\) factor appears on both sides because the left outgoing
Rayleigh factor has \(\partial_x=-\ii\gamma_m\), while the left port outward
normal is \(-e_x\).

The port matching equations are
\[
  D_h^\pm+D_s^\pm
  =
  \Din^\pm p_{\mathrm{in}}^\pm+\Dout^\pm r^\pm,
\]
\[
  N_h^\pm+N_s^\pm
  =
  \Nin^\pm p_{\mathrm{in}}^\pm+\Nout^\pm r^\pm .
\]
With
\[
  x=
  \begin{bmatrix}
  \eta\\
  a\\
  b\\
  r^-\\
  r^+
  \end{bmatrix},
\]
the coupled matrix is
\[
  \Atr
  =
  \begin{bmatrix}
  \Aqp & H_\Sigma^a & H_\Sigma^b & 0 & 0\\
  F_- & I & I & -\Dout^- & 0\\
  \ii\GammaR F_- & -\ii\GammaR & \ii\GammaR & -\Nout^- & 0\\
  F_+ & E & E^{-1} & 0 & -\Dout^+\\
  \ii\GammaR F_+ & \ii\GammaR E & -\ii\GammaR E^{-1} & 0 & -\Nout^+
  \end{bmatrix}.
\]
For forced lead scattering,
\[
  \Atr x=b_{\mathrm{in}},
\]
where
\[
  b_{\mathrm{in}}
  =
  \begin{bmatrix}
  0\\
  \Din^-p_{\mathrm{in}}^-\\
  \Nin^-p_{\mathrm{in}}^-\\
  \Din^+p_{\mathrm{in}}^+\\
  \Nin^+p_{\mathrm{in}}^+
  \end{bmatrix}.
\]
For the eigenvalue or TEP-type workflow, one uses \(\Atr x=0\).

\subsection{Dimension count and center homogeneous coefficients}

Let \(N=N_{\mathrm{tot}}\) be the number of boundary quadrature nodes on
\(\Sigma\).  Then \(\eta\in\CC^{2N}\).  The BIE/transmission row contributes
\(2N\) equations.  Each port contributes \(K\) Dirichlet equations and \(K\)
Neumann equations, so the two ports contribute \(4K\) equations.  The total
number of equations is
\[
  2N+4K .
\]
The unknowns are
\[
  \eta\in\CC^{2N},
  \qquad
  a,b\in\CC^K,
  \qquad
  r^\pm\in\CC^{K_{\mathrm{out}}^\pm}.
\]
The total number of unknowns is
\[
  2N+2K+K_{\mathrm{out}}^-+K_{\mathrm{out}}^+ .
\]
The matrix is square when
\[
  K_{\mathrm{out}}^-+K_{\mathrm{out}}^+=2K .
\]
In a non-degenerate lossless second-order setting, one expects
\[
  K_{\mathrm{out}}^\pm\approx K .
\]
This is a diagnostic, not an unconditional theorem.  Deviations may indicate
unit-circle propagating modes, band edges, grazing Rayleigh orders, mode
collisions, thresholding issues, or insufficient Rayleigh truncation.

The center homogeneous coefficients \(a,b\) are required.  The density
\(\eta\) represents the scattered correction generated by the center inclusion;
it does not span the full homogeneous solution space in the center cell.  If
\(a,b\) are omitted, the port matching system is generally overdetermined and
does not represent the intended center-cell solution class.
